\documentclass[a4paper,11pt]{article}
\usepackage[utf8]{inputenc}
\usepackage[french]{babel}
\usepackage[T1]{fontenc}
\usepackage{amsmath,amssymb,amsthm}
\usepackage{geometry}
\geometry{margin=2cm}
\usepackage{enumitem}
\usepackage{hyperref}
\usepackage{booktabs}
\usepackage{array}
\usepackage{xcolor}
\usepackage{listings}
\usepackage{titlesec}
\usepackage{fancyhdr}
\usepackage{lastpage}
\usepackage{graphicx}
\usepackage{etoolbox}
\AtBeginDocument{\shorthandoff{;:!?}}

\titleformat{\section}{\large\bfseries\color{blue!60!black}}{}{0pt}{}[\vspace{-0.5ex}\rule{\textwidth}{0.4pt}]
\titleformat{\subsection}{\bfseries\color{blue!40!black}}{}{0pt}{}
\titleformat{\subsubsection}{\itshape\color{blue!30!black}}{}{0pt}{}

\lstdefinestyle{pystyle}{
  frame=single,
  numbers=left,
  numbersep=5pt,
  breaklines=true,
  basicstyle=\small\ttfamily,
  backgroundcolor=\color{gray!5},
  rulecolor=\color{gray!40},
  columns=fullflexible,
  keepspaces=true,
  showstringspaces=false,
  language=Python
}
\lstset{style=pystyle}

\pagestyle{fancy}
\fancyhf{}
\fancyhead[L]{\small STA211 -- TP Deep Learning avec Keras}
\fancyhead[R]{\small Fiche r\'ecapitulative}
\fancyfoot[C]{\thepage/\pageref{LastPage}}
\renewcommand{\headrulewidth}{0.4pt}

\theoremstyle{definition}
\newtheorem*{definition}{D\'efinition}
\theoremstyle{remark}
\newtheorem*{remark}{Remarque}
\newtheorem*{important}{\`A retenir}

\title{\textbf{Fiche r\'ecapitulative -- STA211}\\[4pt]
\large Travaux Pratiques \\[2pt]
\small Deep Learning avec Keras et Manifold Untangling}
\author{Nicolas Thome (CNAM)}
\date{14 juin 2026}

\begin{document}
\maketitle
\thispagestyle{empty}

\begin{abstract}
Ce TP pratique la biblioth\`eque Keras sur MNIST~: r\'egression
logistique, MLP, CNN (type LeNet), puis visualise le \emph{manifold
untangling} par t-SNE, enveloppes convexes, ellipses et
Neighborhood Hit.
\end{abstract}

\vfill
\tableofcontents
\newpage

% ===================================================================
\section{Partie 1 : Deep Learning avec Keras}
% ===================================================================

\subsection{Exercice 0 : Chargement des donn\'ees MNIST}

MNIST~: 60\,000 images train, 10\,000 test, $28\times28$,
10 classes (chiffres 0--9).

\begin{lstlisting}
from keras.datasets import mnist
(X_train, y_train), (X_test, y_test) = mnist.load_data()
X_train = X_train.reshape(60000, 784).astype('float32') / 255
X_test  = X_test.reshape(10000, 784).astype('float32') / 255
\end{lstlisting}

Afficher les 200 premi\`eres images avec Matplotlib.

\paragraph{Question~:} espace des images~? $\mathbb{R}^{784}$.

\subsection{Exercice 1 : R\'egression logistique}

\paragraph{Mod\`ele~:} $p(\hat{y}_{c,i}|\mathbf{x}_i) =
\frac{e^{\langle\mathbf{x}_i;\mathbf{w}_c\rangle + b_c}}
{\sum_{c'} e^{\langle\mathbf{x}_i;\mathbf{w}_{c'}\rangle + b_{c'}}}$

\paragraph{Nombre de param\`etres~:} $784\times 10 + 10 = 7\,850$.

\begin{lstlisting}
model = Sequential()
model.add(Dense(10, input_dim=784, name='fc1'))
model.add(Activation('softmax'))
\end{lstlisting}

\paragraph{One-hot encoding~:}
\begin{lstlisting}
from keras.utils import np_utils
Y_train = np_utils.to_categorical(y_train, 10)
\end{lstlisting}

\paragraph{Entropie crois\'ee~:}
$\mathcal{L}_{\mathbf{W},\mathbf{b}}(\mathcal{D}) =
-\frac{1}{N}\sum_{i=1}^N \log(\hat{y}_{c^*,i})$

\paragraph{Convexit\'e~:} oui, LR + cross-entropy = probl\`eme
convexe $\Rightarrow$ minimum global assur\'e avec bon pas.

\begin{lstlisting}
sgd = SGD(learning_rate=0.5)
model.compile(loss='categorical_crossentropy',
              optimizer=sgd, metrics=['accuracy'])
model.fit(X_train, Y_train, batch_size=100,
          epochs=10, verbose=1)
scores = model.evaluate(X_test, Y_test, verbose=0)
\end{lstlisting}

\paragraph{R\'esultat attendu~:} $\sim$92\,\% accuracy test.

\subsection{Exercice 2 : Perceptron multi-couches (MLP)}

Une couche cach\'ee $L=100$, sigmo\"ide, softmax en sortie.

\paragraph{Nombre de param\`etres~:}
$784\times 100 + 100 + 100\times 10 + 10 = 79\,510$.

\begin{lstlisting}
model = Sequential()
model.add(Dense(100, input_dim=784, name='fc1'))
model.add(Activation('sigmoid'))
model.add(Dense(10, name='fc2'))
model.add(Activation('softmax'))
\end{lstlisting}

\paragraph{Non-convexit\'e~:} MLP + sigmo\"ide = probl\`eme non
convexe. Pas de garantie de minimum global.

\paragraph{R\'esultat attendu~:} $\sim$98\,\% accuracy test.

\subsection{Exercice 3 : R\'eseau de neurones convolutif (CNN)}

Reformatage~: $28\times28\times1$.

\begin{lstlisting}
X_train = X_train.reshape(-1, 28, 28, 1)
X_test  = X_test.reshape(-1, 28, 28, 1)
\end{lstlisting}

\paragraph{Architecture LeNet-like~:}
\begin{enumerate}
  \item Conv2D(16, $5\times5$, ReLU) + MaxPooling($2\times2$)
  \item Conv2D(32, $5\times5$, ReLU) + MaxPooling($2\times2$)
  \item Flatten + Dense(100, sigmo\"ide) + Dense(10, softmax)
\end{enumerate}

\begin{lstlisting}
model = Sequential()
model.add(Conv2D(16, kernel_size=(5,5),
          activation='relu', input_shape=(28,28,1)))
model.add(MaxPooling2D(pool_size=(2,2)))
model.add(Conv2D(32, (5,5), activation='relu'))
model.add(MaxPooling2D(pool_size=(2,2)))
model.add(Flatten())
model.add(Dense(100, activation='sigmoid'))
model.add(Dense(10, activation='softmax'))
\end{lstlisting}

\paragraph{R\'esultat attendu~:} $\sim$99\,\% accuracy test.

\paragraph{Sauvegarde du mod\`ele~:}
\begin{lstlisting}
def saveModel(model, savename):
    model_yaml = model.to_yaml()
    with open(savename+".yaml", "w") as f:
        f.write(model_yaml)
    model.save_weights(savename+".h5")
\end{lstlisting}

% ===================================================================
\section{Partie 2 : Deep Learning et Manifold Untangling}
% ===================================================================

\subsection{Exercice 4 : Visualisation avec t-SNE}

\paragraph{t-SNE~:} r\'eduction non lin\'eaire en 2D qui pr\'eserve
les proximit\'es entre points.

\begin{lstlisting}
from sklearn.manifold import TSNE
tsne = TSNE(n_components=2, init='pca',
            perplexity=30, verbose=2)
X_r = tsne.fit_transform(X_test[:1000])
\end{lstlisting}

\paragraph{Trois m\'etriques de s\'eparation~:}
\begin{enumerate}
  \item \textbf{Enveloppe convexe (ConvexHull)}~: p\'erim\`etre
    minimum contenant tous les points d'une classe.
  \item \textbf{Ellipse de meilleure approximation}~: ajustement
    d'une gaussienne 2D par classe (GaussianMixture).
  \item \textbf{Neighborhood Hit (NH)}~: pour chaque point,
    proportion des $k$ plus proches voisins de la m\^{e}me classe.
    Moyenn\'ee sur toute la base.
\end{enumerate}

\begin{lstlisting}
def neighboring_hit(points, labels, k=6):
    nbrs = NearestNeighbors(n_neighbors=k+1,
                            algorithm='ball_tree').fit(points)
    distances, indices = nbrs.kneighbors(points)
    txs = 0.0
    for i in range(len(points)):
        tx = sum(1 for j in range(1, k+1)
                 if labels[indices[i,j]] == labels[i]) / k
        txs += tx
    return txs / len(points)
\end{lstlisting}

\paragraph{Comparaison ACP vs t-SNE~:} t-SNE s\'epare mieux
les classes non lin\'eairement s\'eparables.

\subsection{Exercice 5 : Repr\'esentations internes des r\'eseaux}

\paragraph{Objectif~:} visualiser l'effet de \emph{manifold
untangling} -- les r\'eseaux de neurones transforment un espace
d'entr\'ee peu s\'eparable en un espace latent lin\'eairement
s\'eparable.

\paragraph{Chargement du mod\`ele~:}
\begin{lstlisting}
def loadModel(savename):
    with open(savename+".yaml", "r") as f:
        model = model_from_yaml(f.read())
    model.load_weights(savename+".h5")
    return model
\end{lstlisting}

\paragraph{Extraction des repr\'esentations latentes~:}
\begin{lstlisting}
model.pop()  # supprime la couche softmax
model.pop()  # supprime la couche FC
representations = model.predict(X_test)
\end{lstlisting}

\paragraph{R\'esultat~:} Les repr\'esentations internes (MLP cach\'e
100D ou CNN) montrent une bien meilleure s\'eparation des classes
que l'espace d'entr\'ee original, visible par t-SNE.

\begin{important}
  Le \emph{manifold untangling} est la capacit\'e des r\'eseaux
  profonds \`a \guillemotleft{}d\'em\^eler\guillemotright{} les
  vari\'et\'es des donn\'ees, rendant les classes lin\'eairement
  s\'eparables dans l'espace des repr\'esentations internes.
\end{important}

% ===================================================================
\newpage
\section*{Questions cl\'es (QCM)}
% ===================================================================
\addcontentsline{toc}{section}{Questions cl\'es (QCM)}

\begin{enumerate}
  \item Combien de param\`etres un mod\`ele de r\'egression
    logistique a-t-il sur MNIST ($28\times28$, 10 classes)~?
    \begin{enumerate}
      \item 784
      \item 7\,850
      \item 10\,000
      \item 79\,510
    \end{enumerate}
    \textbf{R\'eponse~: b ($784\times10 + 10 = 7\,850$).}

  \item La fonction de co\^ut de la r\'egression logistique
    (cross-entropy) est-elle convexe~?
    \begin{enumerate}
      \item Oui, minimum global garanti
      \item Non, minimum local seulement
      \item Seulement avec SGD
      \item Cela d\'epend de l'initialisation
    \end{enumerate}
    \textbf{R\'eponse~: a (LR + cross-entropy = convexe).}

  \item Combien de param\`etres un MLP avec 1 couche cach\'ee
    de 100 neurones a-t-il sur MNIST~?
    \begin{enumerate}
      \item 7\,850
      \item 79\,510
      \item 100\,000
      \item 797\,000
    \end{enumerate}
    \textbf{R\'eponse~: b ($784\times100{+}100{+}100\times10{+}10$).}

  \item Quelle accuracy peut-on attendre du MLP sur MNIST~?
    \begin{enumerate}
      \item 85\,\%
      \item 92\,\%
      \item 98\,\%
      \item 99\,\%
    \end{enumerate}
    \textbf{R\'eponse~: c ($\sim$98\,\%).}

  \item Quelle accuracy peut-on attendre du CNN (LeNet-like)
    sur MNIST~?
    \begin{enumerate}
      \item 85\,\%
      \item 92\,\%
      \item 98\,\%
      \item 99\,\%
    \end{enumerate}
    \textbf{R\'eponse~: d ($\sim$99\,\%).}

  \item Que signifie l'acronyme t-SNE~?
    \begin{enumerate}
      \item \emph{t-Statistics and Neural Embedding}
      \item \emph{t-Distributed Stochastic Neighbor Embedding}
      \item \emph{Topological Sample Neighborhood Estimation}
      \item \emph{Tensor-Structured Nonlinear Embedding}
    \end{enumerate}
    \textbf{R\'eponse~: b.}

  \item Que mesure la m\'etrique Neighborhood Hit (NH)~?
    \begin{enumerate}
      \item La distance moyenne entre classes
      \item Le taux de voisins de m\^{e}me classe parmi les $k$ plus
        proches voisins
      \item La variance intra-classe
      \item Le nombre de clusters
    \end{enumerate}
    \textbf{R\'eponse~: b.}

  \item Quel est le principe du \emph{manifold untangling}~?
    \begin{enumerate}
      \item R\'eduire la dimension des donn\'ees
      \item D\'em\^{e}ler les vari\'et\'es pour rendre les classes
        s\'eparables dans l'espace latent
      \item Augmenter le nombre de param\`etres
      \item Initialiser les poids al\'eatoirement
    \end{enumerate}
    \textbf{R\'eponse~: b.}

  \item Comment extraire la couche cach\'ee d'un mod\`ele Keras~?
    \begin{enumerate}
      \item model.predict(layer=hidden)
      \item model.pop() puis model.predict()
      \item model.get\_layer('hidden').output
      \item model.forward(hidden=True)
    \end{enumerate}
    \textbf{R\'eponse~: b ou c (pop + predict ou get\_layer).}

  \item Quelle m\'ethode de visualisation non lin\'eaire est
    utilis\'ee dans le TP~?
    \begin{enumerate}
      \item ACP
      \item t-SNE
      \item LDA
      \item k-means
    \end{enumerate}
    \textbf{R\'eponse~: b (t-SNE).}

\end{enumerate}

% ===================================================================
\begin{thebibliography}{9}
% ===================================================================

\bibitem{chollet2015}
  F.~Chollet.
  \newblock \emph{Keras}. \url{https://keras.io}, 2015.

\bibitem{tsne08}
  L.~van der Maaten, G.~Hinton.
  \newblock \emph{Visualizing Data using t-SNE}.
  \newblock JMLR, 9:2579--2605, 2008.

\bibitem{lecun1989}
  Y.~LeCun et al.
  \newblock \emph{Backpropagation applied to handwritten zip code
    recognition}.
  \newblock Neural Computation, 1(4):541--551, 1989.

\end{thebibliography}

\end{document}
